\section*{Ex.~2.2 (5 pts) --- Why process creation by a running process is the general case}

The first two cases describe what \emph{prompts} a process to be created; they do
not describe a different creator. They are not alternatives to the third case, they
are the third case seen from the user's side.

When someone double-clicks an icon or types a command, that person never touches the
kernel. The click and the keystrokes are input events, and the kernel delivers them
to a process that is \emph{already running}: the desktop shell or window manager in
the first instance, the command shell in the second. That process interprets the
event, decides which program was meant, and calls \texttt{fork()} itself. The shell
is the creator; the human only supplied the reason. Case~2 differs from case~1 only
in where the trigger originated (another process, a timer, an arriving network
connection) while the mechanism, a running process calling \texttt{fork()}, is
unchanged.

All three cases therefore share a single creation primitive and differ only in what
caused the running process to invoke it. Unix makes this literal: exactly one process
is not born this way, the initial process (\texttt{init} or \texttt{systemd}, pid~1)
that the kernel hand-builds at boot. Every other process is the child of some running
process.

\section*{Ex.~2.3 (5 pts) --- A program using \texttt{fork()} and \texttt{exec()} to run \texttt{ls}}

\lstdefinestyle{p2cstyle}{
  language=C,
  basicstyle=\ttfamily\small,
  keywordstyle=\color{blue!60!black},
  commentstyle=\itshape\color{gray!50!black},
  stringstyle=\color{red!50!black},
  numbers=left,
  numberstyle=\tiny\color{gray},
  numbersep=8pt,
  showstringspaces=false,
  breaklines=true,
  frame=single,
  framesep=4pt,
  columns=fullflexible,
  keepspaces=true,
}

\begin{lstlisting}[style=p2cstyle]
#include <stdio.h>
#include <stdlib.h>
#include <unistd.h>
#include <sys/wait.h>

int main(void)
{
    pid_t pid = fork();

    if (pid < 0) {                      /* case 1: fork failed */
        perror("fork");
        return EXIT_FAILURE;
    }

    if (pid == 0) {                     /* case 2: we are the child */
        char *argv[] = { "ls", "-l", NULL };
        execvp("ls", argv);
        /* only reached if execvp failed: the image was never replaced */
        perror("execvp");
        _exit(127);
    }

    /* case 3: we are the parent, pid holds the child's pid */
    int status;
    if (waitpid(pid, &status, 0) < 0) {
        perror("waitpid");
        return EXIT_FAILURE;
    }

    if (WIFEXITED(status))
        printf("parent: child %d exited with status %d\n",
               (int) pid, WEXITSTATUS(status));
    else if (WIFSIGNALED(status))
        printf("parent: child %d killed by signal %d\n",
               (int) pid, WTERMSIG(status));

    return EXIT_SUCCESS;
}
\end{lstlisting}

\texttt{fork()} returns twice, so its return value is the only thing that tells the
two processes apart, and all three of its outcomes are handled. A negative value
means no child was created. In the child (return value $0$) \texttt{execvp} replaces
the process image with \texttt{/usr/bin/ls} while keeping the same pid, searching
\texttt{PATH} for the binary and passing a \texttt{NULL}-terminated \texttt{argv}. A
successful \texttt{exec} never returns, so the lines after it run only on failure,
where they report the error and leave via \texttt{\_exit} rather than \texttt{exit}
so the child does not flush the parent's copy of the \texttt{stdio} buffers. In the
parent (return value the child's pid) \texttt{waitpid} blocks until that specific
child terminates, which both sequences the output and reaps the child so it does not
linger as a zombie.

Compiled with \texttt{gcc -Wall -Wextra -o forkexec forkexec.c} (gcc 12.2.0, Debian
Linux) and run in a directory holding two small files. The actual output was:

\begin{lstlisting}[basicstyle=\ttfamily\small,frame=single,framesep=4pt,columns=fullflexible]
total 8
-rw-r--r-- 1 xing xing 1 Sep 21 16:37 alpha.txt
-rw-r--r-- 1 xing xing 2 Sep 21 16:37 beta.txt
parent: child 100 exited with status 0
\end{lstlisting}

The first three lines come from \texttt{ls} after it has taken over the child, the
last from the parent after \texttt{waitpid} returns. Setting \texttt{PATH} to a
directory that does not exist exercises the failure branch: \texttt{execvp: No such
file or directory}, then \texttt{parent: child 107 exited with status 127}.
